% GENERATED, edit the French unit or the translation table
% GENERATED from content/data/enrichment.json by mxbuild.enrichment; do not edit.
\mxunit{enrichment-waiting-time-solutions}{Waiting for the first success}
\mxsubtitle{Upper high school · probability and sequences \textperiodcentered{} Worked solutions \textperiodcentered{} 08}

\begin{exercise}
With $p=\frac{1}{4}$, find $P(T=3)$.
\end{exercise}
\begin{solution}
$(\frac{3}{4})^{2}(\frac{1}{4})=\frac{9}{64}$.
\end{solution}

\begin{exercise}
With $p=\frac{1}{4}$, find $P(T\leq 3)$.
\end{exercise}
\begin{solution}
$1-P(T>3)=1-(\frac{3}{4})^{3}=\frac{37}{64}$.
\end{solution}

\begin{exercise}
For $0<p<1$, find $P(T>m+n|T>m)$.
\end{exercise}
\begin{solution}
The ratio of tail probabilities is $\frac{(1-p)^{m+n}}{1-p}^{m}=(1-p)^{n}$. After m failures, the remaining waiting-time distribution is unchanged.
\end{solution}

\begin{exercise}
Success has probability $\num{0,2}$ at each independent trial. What is the mean waiting time? Must the fifth trial succeed?
\end{exercise}
\begin{solution}
$E(T)=5$ trials. No: every trial still has success probability $\num{0,2}$, independently of previous failures.
\end{solution}

\begin{problem}[Investigate, justify, conclude]
Roll a fair die until the first 6, but stop after 10 rolls in any case. Let N be the number of rolls actually made. Find $P(N=10)$, the probability of obtaining at least one 6, and E(N). For the mean, use $E(N)=\sum  P(N>n)$, from $n=0$ to 9.
\end{problem}
\begin{solution}
$N=10$ when the first nine rolls fail, regardless of the tenth: $P(N=10)=(\frac{5}{6})^{9}$. The probability of at least one 6 is $1-(\frac{5}{6})^{10}$. For $0\leq n\leq 9$, $P(N>n)=(\frac{5}{6})^{n}$, giving $E(N)=6[1-(\frac{5}{6})^{10}]$, about $\num{5,031}$ rolls. Truncation reduces the expectation below 6.
\end{solution}

